
\documentclass[aspectratio=169]{beamer}

% Theme and appearance
\usetheme{Madrid}
\usecolortheme{default}

% Packages
\usepackage{amsmath}

% Presentation information
\title{How to Learn Mathematics}
\subtitle{Developing Understanding, Reasoning, and Independence}
\author{}
\institute{}
\date{}

\begin{document}

%------------------------------------------------
% Title slide
%------------------------------------------------
\begin{frame}
    \titlepage
\end{frame}

%------------------------------------------------
% Section 1: What is mathematics about?
%------------------------------------------------
\begin{frame}{What Is Mathematics About?}

    \begin{itemize}
        \item Mathematics is more than memorizing formulas
              and procedures.
        \pause

        \item It is fundamentally about
              \textbf{logic, reasoning, and understanding}.
        \pause

        \item The goal is not just to know \emph{how} to solve
              a problem, but also to understand \emph{why}
              a method works.
        \pause

        \item Memorization has its place, but it should support
              understanding, not replace it.
        \pause
    \end{itemize}

\end{frame}

%------------------------------------------------
% Section 2A: Explain concepts to others
%------------------------------------------------
\begin{frame}{Two Important Ways to Learn Mathematics}

    \begin{block}{Method 1: Explain Concepts to Others}
        \begin{itemize}
            \item Explaining an idea forces you to organize
                  your thoughts.
            \pause

            \item If you truly understand a concept, you should
                  be able to explain it in your own words.
            \pause

            \item Don't just explain \emph{what} to do.
                  Explain \emph{why} it works.
            \pause

            \item Teaching others is also a way of testing
                  your own understanding.
            \pause
        \end{itemize}
    \end{block}

\end{frame}

%------------------------------------------------
% Section 2B: Analyze your mistakes
%------------------------------------------------
\begin{frame}{Two Important Ways to Learn Mathematics}

    \begin{block}{Method 2: Analyze Your Mistakes}
        \begin{itemize}
            \item Mistakes are opportunities to learn,
                  not merely failures.
            \pause

            \item Don't just correct a wrong answer.
                  Find out \textbf{why} you got it wrong.
            \pause

            \item Identify the first incorrect step and
                  understand the misconception behind it.
            \pause

            \item Learn from the mistake so you can avoid
                  repeating it.
            \pause
        \end{itemize}
    \end{block}

\end{frame}

%------------------------------------------------
% Section 3: The bottom line
%------------------------------------------------
\begin{frame}{The Bottom Line}

    \begin{itemize}
        \item Don't just memorize. \textbf{Understand.}
        \pause

        \item Don't just follow steps. \textbf{Reason.}
        \pause

        \item Don't just correct mistakes.
              \textbf{Learn from them.}
        \pause

        \item \textbf{Take ownership of your learning and
              become an independent mathematical thinker.}
        \pause
    \end{itemize}

\end{frame}

%------------------------------------------------
% Closing slide
%------------------------------------------------
\begin{frame}[plain]
    \centering
    \vfill

    {\Huge\bfseries Learn to Think Mathematically}

    \vspace{1cm}

    {\large Understand. Explain. Analyze. Improve.}

    \vfill
\end{frame}

\end{document}
